% Unnumbered diagnostic only; no claimed hereditary theorem or counterexample to it.
\subsection*{A forced nonvanishing on a one-root rational bridge}

This calculation tests the proposed componentwise pole-stripping condition
\[
 H^0\bigl(R,\mathcal End(W)\otimes L\bigr)=0
       \quad\text{whenever }\deg L<0.
\]
It shows that this condition is false even for the forced, balanced inertia
of a minimal-rank generating anchor on a bridge. It does not show that a
globally compatible pole section exists, or that another stripping order
cannot work.

Let $m>1$ be prime to the characteristic of the algebraically closed field
$k$, and let $\pi:R=\sqrt[m]{(\mathbf P^1,\infty)}\to\mathbf P^1$.
Write $D$ for the root divisor, so $\mathcal O_R(mD)=\pi^*\mathcal O(1)$.
Let $W$ be any rank-$m$ vector bundle whose inertia representation at $D$
contains each character exactly once. No splitting type, deformation
normalization or canonical-order stability is assumed. Put $V=\pi_*W$.
Then $V$ is a rank-$m$ vector bundle, and
\[
 \chi\bigl(R,\mathcal End(W)(-D)\bigr)=\frac{m(m-1)}2>0.
                                                        \qquad(*)
\]
In particular $H^0(R,\mathcal End(W)(-D))\ne0$, although
$\deg\mathcal O_R(-D)=-1/m$.

Here is a direct proof, avoiding any orbifold Riemann--Roch convention.
The natural map
\[
 \pi_*\mathcal End(W)(-D)\longrightarrow\mathcal End(V)
\]
is injective and is an isomorphism off $\infty$. On the completed root
chart $k[[z]]$, with $t=z^m$, choose homogeneous frames and order them so
that invariant generators of $V$ are $z^{r_i}e_i$, where
$0=r_1<r_2<\cdots<r_m=m-1$. An entry of
$\mathcal End(W)(-D)$ is divisible by $z$. After passing to these invariant
frames, the $(i,j)$ matrix entry is arbitrary over $k[[t]]$ if $r_i>r_j$,
and lies in $t k[[t]]$ if $r_i\le r_j$.
Indeed its smallest allowed exponent $k\ge1$ satisfies
$k\equiv r_i-r_j\pmod m$, and its exponent in the invariant frames is
$r_j+k-r_i$, equal to zero or $m$ in those two cases.
The cokernel is therefore a skyscraper of length $m(m+1)/2$.
Since $\deg\mathcal End(V)=0$, ordinary Riemann--Roch on $\mathbf P^1$
gives $\chi(\mathcal End(V))=m^2$. Exactness of tame coarse pushforward
now gives $(*)$. The calculation is valid for every such $W$, not only
a split example.

\paragraph{Occurrence in the refined residual geometry.}
Consider a new exceptional bridge obtained by blowing up a node of root
orders $(1,m)$ with zero secondary residue. Its own root order is one,
its self-intersection is $-1$, and its spectral curve is the above
one-root $\mathbf P^1$. A rank-$m$ generating anchor has precisely the
character multiplicities stipulated above. Let $E$ denote the bridge and
let $D_{\mathrm{right}}$ denote the root divisor of the adjacent
order-$m$ component. The effective Cartier divisor on the root surface
\[
                         E+(m-1)D_{\mathrm{right}}
\]
restricts to
$\mathcal O_{\mathbf P^1}(-1)\otimes\mathcal O_R((m-1)D)
=\mathcal O_R(-D)$ on the bridge. Its rational intersection degree there
is negative, yet the displayed endomorphism-section space is nonzero.
Thus a rule requiring the proposed vanishing on every negatively
intersected component cannot be inferred from minimal rank, regular
inertia multiplicities, or reduced spectral balance.

\paragraph{Exact limitations.}
This is a residual-component calculation. A nonzero section here need
not extend over the whole exceptional cycle or match either neighboring
component. The effective divisor above may have another negatively
intersected component on which a stripping step is possible. Nor does the
calculation assert the existence of a bad globally canonical full order,
an unbounded cohomology family, or a failure of the desired hereditary
count. A proof by coupled whole-cycle cancellation, or by a suitable
choice of stripping order, remains possible. The calculation merely
excludes the stated universal componentwise vanishing shortcut.

\paragraph{Source scope.}
All new assertions are proved by the displayed homogeneous matrix
calculation and ordinary vector-bundle Riemann--Roch on $\mathbf P^1$.
The already audited tame root-chart description and exact coarse
invariants are reused. This note is not a numbered or promoted result
and has not received four independent audits.
